% !TeX program = lualatex
% =============================================================================
% bingo template — math-bingo cards
% =============================================================================
%
% Requires LuaLaTeX (the randomizer runs in Lua):
%   lualatex bingo-template.tex

% All options are class options. Card identity (course-*, institution,
% bingo-title) fills the header/footer; without it the slots show placeholders or
% are populated from a coursemeta.tex up the directory tree. show-instructions
% prints the how-to-play block on every card (default wording in
% bingo-instructions.tex; override with bingo-instructions={...} for boxed inline
% text or bingo-instructions-file=my-notes for an \input'd file).
\documentclass[
	course-subject    = Math,
	course-number     = 181,
	course-title      = {Calculus I},
	institution       = {University of Nevada, Reno},
	bingo-title       = {Exam 3 Bingo},
	show-instructions = true,
]{bingo}

% The answer pool — a comma-separated list, one entry per line. Cells are math
% mode automatically (no $ needed). The list may hold more than 25 entries; \free
% is the free space. Brace an entry that itself contains a top-level comma.
\bingobank{
	xe^{-x},
	\pi,
	\alpha,
	24x-12,
	e^x,
	\tfrac{5x}{e^{5x}},
	4\sqrt{5},
	\free,
	x\ln x,
	\beta,
	7xe^x(2+x),
	\sin x,
	x^3,
	\delta,
	\tfrac{1}{x},
	\frac{\ln(e^x+x)}{x},
	\frac{x^2}{\sin(2x)},
	x+\frac{7500}{x},
	\frac{e^x+1}{e^x+x},
	\frac{18}{\sqrt[3]{9}},
	\frac{15}{4},
	\epsilon,
	2\pi r^2+\frac{324\pi}{r},
	12x^2(x-1),
	\ln x-1,
	\gamma,
	\sqrt{\frac{13}{3}},
	3\sqrt[3]{3},
	x^2,
	50\sqrt{3}
}

\begin{document}

% With show-instructions=true (above), the how-to-play block prints on every
% card below. For a single instructions cover page instead, drop that option and
% call \bingoinstructions here followed by \newpage.

% A one-off card: place entries by hand in a 5x5 grid (& between cells, \\
% between rows). Cells are math mode automatically.
\begin{bingocard}
	\beta & \pm & \Omega & \mu & \mathbb{R} \\
	\to & \int & \neq & \notin & \psi \\
	\exists & \forall & \free & \frac{d}{dx} & \sqrt{x} \\
	\theta & \circ & \pi & \prod & \gamma \\
	\alpha & \infty & \sigma & \sum & \Delta \\
\end{bingocard}

\newpage

% A randomized set from the pool — one card per page. Fixed across builds; set
% `seed` to deal a fresh set. `keepfree` (default) guarantees one free space.
\bingocards[copies=3, randomize]

\end{document}

% =============================================================================
% More options (see README.md):
%
%	% reshuffle: cards are fixed across builds; set seed to deal a new set:
%	\bingocards[copies=30, randomize, seed=2]
%
%	% free space as an ordinary pool entry (a large pool may then drop it):
%	\bingocards[copies=30, randomize, keepfree=false]
%
%	% a second named pool, emitted by name:
%	\bingobank[final]{ ... }
%	\bingocards[copies=23, randomize, bank=final]
%
%	% more class options (add to the \documentclass list above):
%	%   cell-size=3.0cm, free-symbol=\cdot, show-headers=false, show-title=true
%	%   bingo-instructions={Mark answers as they are read; five in a row wins.}
%	%   bingo-instructions-file=my-instructions   % unboxed file, full control
% =============================================================================
